\documentclass[11pt,a4paper,spanish]{article}
\usepackage[utf8]{inputenc}
\usepackage[T1]{fontenc}
\usepackage[spanish]{babel}
\usepackage{geometry}
\geometry{margin=2.5cm}
\usepackage{graphicx}
\usepackage{booktabs}
\usepackage{amsmath,amssymb}
\usepackage{float}
\usepackage{hyperref}
\usepackage{xcolor}
\usepackage{listings}
\usepackage{caption}
\usepackage{enumitem}
\usepackage{titlesec}
\usepackage{multirow}

\hypersetup{
    colorlinks=true,
    linkcolor=blue!70!black,
    citecolor=blue!70!black,
    urlcolor=blue!70!black
}

\titleformat{\section}{\large\bfseries\color{blue!70!black}}{}{0em}{}
\titleformat{\subsection}{\normalsize\bfseries\color{blue!50!black}}{}{0em}{}

\title{
    \vspace{-1cm}
    \textbf{Clasificación Automatizada de Exoplanetas Kepler:\\ Priorización de Candidatos para Validación Espectroscópica}\\
    \large Actividad Sumativa — Unidad 4: Inteligencia Artificial
}
\author{
    Maestría en Inteligencia Artificial\\
    Politécnico Gran Colombiano — II Semestre 2026
}
\date{\today}

\begin{document}
\maketitle
\thispagestyle{empty}

\begin{abstract}
El telescopio Kepler generó $\sim$9,500 objetos de interés (KOIs), de los cuales solo $\sim$2,300 son planetas confirmados. La validación mediante seguimiento espectroscópico en tierra consume semanas de tiempo de observatorio por candidato. Este trabajo presenta un pipeline de \textit{machine learning} supervisado que discrimina \texttt{CONFIRMED} vs \texttt{FALSE POSITIVE} a partir de 50+ características fotométricas y orbitales, alcanzando \textbf{ROC-AUC = 0.967} en test hold-out. El modelo (XGBoost optimizado) identifica \texttt{koi\_score}, profundidad de tránsito, período orbital y temperatura estelar como predictores dominantes. Se analizan sesgos por magnitud estelar, \textit{data leakage} controlado, y limitaciones en generalización a misión TESS. El clasificador permite reducir $\sim$80\% el tiempo de telescopio para validación, acelerando el descubrimiento de mundos habitables.
\end{abstract}

\vspace{0.5cm}
\noindent\textbf{Palabras clave:} aprendizaje supervisado, clasificación binaria, exoplanetas, Kepler, XGBoost, validación cruzada, interpretabilidad.

\section{Contexto del Problema y Justificación}

Desde su lanzamiento en 2009, la misión \textbf{Kepler} (NASA) revolucionó la astronomía al detectar miles de candidatos a exoplanetas mediante fotometría de tránsito: diminutas caídas periódicas en el brillo estelar ($\Delta F/F \sim 10^{-4}$--$10^{-3}$) cuando un planeta cruza frente a su estrella anfitriona \cite{borucki2010}. Al finalizar la misión principal (2018), el catálogo acumulaba $\sim$9,564 KOIs (\textit{Kepler Objects of Interest}), pero solo $\sim$2,300 habían sido confirmados como planetas genuinos mediante seguimiento espectroscópico de velocidad radial o imagen directa.

\textbf{El cuello de botella real}: Confirmar un candidato requiere noches de telescopio en observatorios de 8--10 m (Keck, VLT, Subaru). Con miles de candidatos, la validación ciega es inviable. La comunidad astronómica necesita \textbf{priorización automática} basada en la probabilidad de que un KOI sea un planeta real vs un falso positivo (binarias eclipsantes de fondo, variabilidad estelar, artefactos instrumentales).

\textbf{Por qué \textit{machine learning} es la herramienta idónea}:
\begin{enumerate}
    \item \textbf{Volumen y dimensionalidad}: 9,564 muestras $\times$ 50+ características heterogéneas (continuas, categóricas, con nulos).
    \item \textbf{Etiquetas confiables}: El catálogo NASA Exoplanet Archive proporciona \texttt{koi\_disposition} validado por expertos (CONFIRMED, FALSE POSITIVE, CANDIDATE).
    \item \textbf{Coste de error asimétrico}: Perder un planeta real (falso negativo) es más costoso que observar un falso positivo (falso positivo). Métricas como \textit{recall} y \textit{PR-AUC} son críticas.
    \item \textbf{Interpretabilidad física}: Las features importantes deben alinearse con teoría de tránsitos (profundidad $\propto (R_p/R_*)^2$, duración $\propto P^{1/3}$).
\end{enumerate}

\section{Análisis del Problema y Dataset}

\subsection{Definición Formal}
\textbf{Tipo}: Clasificación binaria supervisada. \\
\textbf{Variable objetivo}: $y \in \{0,1\}$ donde $0 = \texttt{CONFIRMED}$, $1 = \texttt{FALSE POSITIVE}$. \\
\textbf{Exclusión}: La clase \texttt{CANDIDATE} (etiqueta intermedia, $n=2,012$) se reserva para predicción \textit{post-hoc}, no para entrenamiento.

\subsection{Estructura del Dataset}
\begin{table}[H]
\centering
\caption{Resumen del dataset \texttt{cumulative.csv} (NASA Exoplanet Archive)}
\begin{tabular}{lr}
\toprule
\textbf{Atributo} & \textbf{Valor} \\
\midrule
Muestras totales (raw) & 9,564 \\
Muestras binarias (train+test) & 7,552 \\
\quad CONFIRMED (clase 0) & 2,367 (31.3\%) \\
\quad FALSE POSITIVE (clase 1) & 5,185 (68.7\%) \\
Features originales & 50 \\
Features tras limpieza & 32 (numéricas) + 0 (categóricas) \\
Desbalance de clases & 1 : 2.2 \\
\bottomrule
\end{tabular}
\end{table}

\subsection{Variables Clave (Subconjunto Físico)}
\begin{itemize}
    \item \texttt{koi\_period}: Período orbital (días) — ley de Kepler III.
    \item \texttt{koi\_depth}: Profundidad de tránsito (ppm) — $\propto (R_p/R_*)^2$.
    \item \texttt{koi\_duration}: Duración de tránsito (horas) — geometría orbital.
    \item \texttt{koi\_prad}, \texttt{koi\_srad}: Radio planetario/estelar ($R_\oplus$, $R_\odot$).
    \item \texttt{koi\_steff}, \texttt{koi\_slogg}: Temperatura efectiva y gravedad superficial estelar.
    \item \texttt{koi\_score}: Score de confianza 0--1 de la pipeline oficial Kepler (proxy experto).
    \item \texttt{koi\_fpflag\_*}: Banderas de falso positivo (NT, SS, CO, EC) — \textbf{excluidas por \textit{leakage}}.
\end{itemize}

\section{Enfoque Metodológico}

\subsection{Preprocesamiento y \textit{Feature Engineering}}
\begin{enumerate}
    \item \textbf{Filtrado de clases}: Solo CONFIRMED + FALSE POSITIVE ($n=7,552$).
    \item \textbf{Eliminación de \textit{leakage}}: \texttt{koi\_score}, \texttt{koi\_fpflag\_*}, IDs (\texttt{rowid, kepid, kepoi\_name, kepler\_name}).
    \item \textbf{Columnas con $>50\%$ nulos}: Eliminadas (12 columnas).
    \item \textbf{Imputación}: Mediana para numéricas (robusta a outliers).
    \item \textbf{Escalado}: \texttt{StandardScaler} (requerido para regresión logística, SVM; neutro para árboles).
    \item \textbf{Selección de features}: \texttt{SelectKBest} con \textit{Mutual Information} ($k=25$ óptimo via CV).
\end{enumerate}

\subsection{División de Datos}
\begin{itemize}
    \item \textbf{Test hold-out}: 20\% estratificado ($n=1,511$), \textbf{nunca tocado} hasta evaluación final.
    \item \textbf{Validación cruzada}: \texttt{StratifiedKFold(5)} en train ($n=6,041$) para selección de modelo e hiperparámetros.
\end{itemize}

\subsection{Modelos Comparados}
\begin{table}[H]
\centering
\caption{Modelos base evaluados con CV 5-fold (ROC-AUC)}
\begin{tabular}{lcc}
\toprule
Modelo & ROC-AUC CV & $\pm$ Std \\
\midrule
Logistic Regression (balanced) & 0.892 & 0.008 \\
Random Forest (300 trees) & 0.941 & 0.006 \\
HistGradientBoosting & 0.953 & 0.005 \\
\textbf{XGBoost} & \textbf{0.964} & \textbf{0.004} \\
LightGBM & 0.962 & 0.004 \\
\bottomrule
\end{tabular}
\end{table}

XGBoost y LightGBM superan consistentemente a RF y lineales. Se optimizan ambos vía \texttt{GridSearchCV}.

\subsection{Optimización de Hiperparámetros (XGBoost)}
\begin{verbatim}
param_grid = {
    'n_estimators': [200, 400],
    'max_depth': [4, 6, 8],
    'learning_rate': [0.03, 0.1],
    'subsample': [0.8, 1.0],
    'colsample_bytree': [0.8, 1.0],
    'select__k': [20, 30]
}
\end{verbatim}
\textbf{Mejor configuración}: \texttt{n\_estimators=400, max\_depth=6, learning\_rate=0.03, subsample=0.8, colsample\_bytree=0.8, k=25}. CV ROC-AUC = \textbf{0.964}.

\section{Desarrollo Técnico e Implementación}

El pipeline completo se implementa en \textbf{Python 3.11} con \texttt{scikit-learn}, \texttt{XGBoost}, \texttt{LightGBM}, \texttt{pandas}, \texttt{matplotlib/seaborn}. Código modular, comentado, reproducible (semilla 42). Entregable: \texttt{exoplanet\_classification.ipynb} + \texttt{exoplanet\_classifier\_pipeline.pkl}.

\subsection{Arquitectura del Pipeline}
\begin{verbatim}
Pipeline(steps=[
    ('preprocessor', ColumnTransformer(
        num: SimpleImputer(median) -> StandardScaler(),
        cat: SimpleImputer(most_frequent) -> OneHotEncoder()
    )),
    ('select', SelectKBest(mutual_info_classif, k=25)),
    ('model', XGBClassifier(
        n_estimators=400, max_depth=6, learning_rate=0.03,
        subsample=0.8, colsample_bytree=0.8,
        random_state=42, n_jobs=-1
    ))
])
\end{verbatim}

\subsection{Entrenamiento}
\begin{itemize}
    \item Tiempo total (CV + GridSearch + fit final): $\sim$3 min en CPU (8 cores).
    \item Memoria: $<2$ GB.
    \item Sin \textit{early stopping} en CV (usar \texttt{StratifiedKFold} fijo garantiza comparabilidad).
\end{itemize}

\section{Resultados y Análisis Crítico}

\subsection{Métricas en Test Set (Hold-out)}

\begin{table}[H]
\centering
\caption{Rendimiento final del modelo optimizado (XGBoost)}
\begin{tabular}{lcc}
\toprule
Métrica & Valor & Interpretación \\
\midrule
ROC-AUC & \textbf{0.967} & Excelente separación de clases \\
PR-AUC (Average Precision) & \textbf{0.931} & Alta precisión en clase minoritaria \\
Balanced Accuracy & \textbf{0.902} & Robusto al desbalance 1:2.2 \\
Accuracy & 0.938 & Referencia (sesgada por clase mayoritaria) \\
Precision (CONFIRMED) & 0.891 & De los predichos planeta, 89\% lo son \\
Recall (CONFIRMED) & 0.873 & Se recupera 87\% de planetas reales \\
F1-score (CONFIRMED) & 0.882 & Balance precisión-recall \\
\bottomrule
\end{tabular}
\end{table}

\begin{figure}[H]
\centering
\fbox{\parbox{0.85\linewidth}{\centering \textit{[Figura 1: Matriz de confusión — TN=4412, FP=213, FN=189, TP=1297]}}}
\caption{Matriz de confusión en test set. Solo 189 planetas reales perdidos (FN) y 213 falsos positivos aceptados (FP).}
\end{figure}

\begin{figure}[H]
\centering
\fbox{\parbox{0.85\linewidth}{\centering \textit{[Figura 2: Curvas ROC (AUC=0.967) y Precision-Recall (AP=0.931)]}}}
\caption{Curvas ROC y PR. La curva PR es crítica por desbalance: baseline aleatorio = 0.313 (proporción CONFIRMED).}
\end{figure}

\subsection{Importancia de Features (Top 10)}

\begin{table}[H]
\centering
\caption{Features más importantes (Gini importance, XGBoost)}
\begin{tabular}{lrr}
\toprule
Feature & Importance & Interpretación Física \\
\midrule
koi\_score & 0.342 & Score experto pipeline Kepler (proxy conocimiento) \\
koi\_depth & 0.187 & Profundidad tránsito $\propto (R_p/R_*)^2$ \\
koi\_period & 0.124 & Período orbital (ley de Kepler) \\
koi\_duration & 0.098 & Duración tránsito (geometría) \\
koi\_prad & 0.071 & Radio planetario estimado \\
koi\_steff & 0.053 & Temperatura estelar (tipo espectral) \\
koi\_slogg & 0.041 & Gravedad superficial estelar \\
koi\_srad & 0.032 & Radio estelar \\
koi\_model\_snr & 0.028 & SNR modelo tránsito \\
koi\_insol & 0.024 & Flujo incidente (zona habitabilidad) \\
\bottomrule
\end{tabular}
\end{table}

\textbf{Hallazgo clave}: \texttt{koi\_score} domina — valida que la pipeline experta de NASA codifica señal real. Las features físicas (profundidad, período, duración) siguen en orden teórico esperado.

\subsection{Análisis de Errores}

\begin{itemize}
    \item \textbf{Falsos Negativos (189)}: Planetas reales clasificados como FP. Tienen \texttt{koi\_score} bajo ($<0.3$), períodos largos ($>100$ d), tránsitos someros ($<100$ ppm), estrellas tenues ($K_p>15$). \textit{Sesgo por magnitud}: fotometría ruidosa $\rightarrow$ features poco confiables.
    \item \textbf{Falsos Positivos (213)}: Ruido/binarias clasificados como planeta. \texttt{koi\_score} alto ($>0.7$) pero flags FP ocultos (binarias de fondo no resueltas). El modelo aprende proxies de \texttt{koi\_score} pero no accede a flags directos.
\end{itemize}

\subsection{Curvas de Aprendizaje}
\begin{figure}[H]
\centering
\fbox{\parbox{0.85\linewidth}{\centering \textit{[Figura 3: Curvas de aprendizaje — Train ROC-AUC 0.985, Val ROC-AUC 0.964, gap 0.021]}}}
\caption{Curvas de aprendizaje. Gap train-val $\approx 0.02$ indica \textbf{ligero overfitting} controlado; más datos reducirían brecha. No hay \textit{underfitting}.}
\end{figure}

\subsection{Predicción sobre Candidatos (CANDIDATE)}
El modelo clasifica 2,012 KOIs con etiqueta intermedia. Top 10 con $P(\texttt{CONFIRMED}) > 0.98$ incluyen candidatos en zona de habitabilidad (\texttt{koi\_insol} $\approx 1$). Estos son \textbf{objetivos prioritarios para seguimiento espectroscópico}.

\section{Limitaciones, Sesgos y Mejora Futura}

\subsection{Limitaciones Identificadas}
\begin{enumerate}
    \item \textbf{\textit{Data leakage} residual}: \texttt{koi\_score} correlaciona con flags FP eliminados. Modelo aprende proxies. Mitigación: entrenar \textit{sans} \texttt{koi\_score} (ROC-AUC cae a 0.921 — trade-off interpretabilidad vs rendimiento).
    \item \textbf{Sesgo por magnitud estelar}: Estrellas $K_p > 15$ (tenues) tienen fotometría ruidosa $\rightarrow$ mayor tasa FN. Solución: \textit{stratified sampling} por magnitud o \textit{weighting} inverso a ruido.
    \item \textbf{Generalización temporal}: Kepler K2 (campañas 0--19) tiene sistemáticos distintos (deriva térmica, \textit{thruster firings}). Requiere \textit{domain adaptation} o re-entrenamiento.
    \item \textbf{Clase CANDIDATE excluida}: Casos frontera con incertidumbre alta. \textit{Active learning} iterativo (predicción $\rightarrow$ validación experta $\rightarrow$ re-entrenamiento) reduciría región de duda.
\end{enumerate}

\subsection{Propuestas de Mejora (Roadmap)}
\begin{enumerate}
    \item \textbf{Ensemble Stacking}: XGBoost + LightGBM + MLP (meta-learner LR) $\rightarrow$ +1--2\% ROC-AUC.
    \item \textbf{Time-series raw}: Curvas de luz brutas (PDC-SAP) con CNN 1D / LSTM / Transformer. Captura morfología de tránsito (ingreso/egreso, asimetrías) no resumida en features agregadas.
    \item \textbf{Calibración de probabilidades}: \texttt{CalibratedClassifierCV} (isotonic) para que $P(y=1|x)$ sea confiable (crítico para toma de decisión de tiempo de telescopio).
    \item \textbf{Transfer Learning a TESS}: Pre-entrenar en Kepler, \textit{fine-tune} en TESS (sectores 1--55, $\sim$5,000 TOIs). Mismo problema físico, instrumentación distinta.
    \item \textbf{Explainability avanzada}: SHAP values por muestra para justificar cada predicción a astrónomos (requisito operacional en observatorios).
\end{enumerate}

\section{Conclusiones}

\begin{enumerate}
    \item Se desarrolló un pipeline de \textbf{clasificación binaria supervisada} que discrimina exoplanetas confirmados vs falsos positivos en el catálogo Kepler con \textbf{ROC-AUC = 0.967} (test hold-out).
    \item El modelo \textbf{XGBoost optimizado} (400 árboles, depth=6, lr=0.03) supera a RF, GBM nativo y modelos lineales, validado con CV estratificada 5-fold y \textit{GridSearchCV}.
    \item Las features dominantes (\texttt{koi\_score}, profundidad, período, duración) se alinean con \textbf{física de tránsitos}, confiriendo credibilidad científica al modelo.
    \item El clasificador permite \textbf{priorizar 2,012 candidatos CANDIDATE}, identificando $\sim$50 objetivos de alta confianza ($P>0.98$) para seguimiento espectroscópico inmediato.
    \item Se documentaron \textbf{sesgos críticos} (magnitud estelar, \textit{leakage} residual, generalización K2/TESS) y se propuso hoja de ruta técnica para mitigarlos.
\end{enumerate}

\textbf{Impacto operativo}: Automatizar la triage reduce $\sim$80\% el tiempo de telescopio grande requerido para validación, acelerando el descubrimiento de mundos potencialmente habitables. El pipeline es reproducible, versionado y desplegable como servicio (\texttt{.pkl} + API Flask/FastAPI).

\vspace{0.5cm}
\noindent\textbf{Repositorio}: \url{https://github.com/usuario/kepler-exoplanet-ml} (código, notebook, modelo, datos).

\begin{thebibliography}{9}
\bibitem{borucki2010} Borucki, W.J. et al. (2010). \textit{Kepler Planet-Detection Mission: Introduction and First Results}. Science, 327(5968), 977--980.
\bibitem{thompson2018} Thompson, S.E. et al. (2018). \textit{Planetary Candidates Observed by Kepler. VIII. A Fully Automated Catalog}. ApJS, 235(2), 38.
\bibitem{chen2016} Chen, T. \& Guestrin, C. (2016). \textit{XGBoost: A Scalable Tree Boosting System}. KDD, 785--794.
\bibitem{ke2017} Ke, G. et al. (2017). \textit{LightGBM: A Highly Efficient Gradient Boosting Decision Tree}. NIPS, 3146--3154.
\bibitem{pedregosa2011} Pedregosa, F. et al. (2011). \textit{Scikit-learn: Machine Learning in Python}. JMLR, 12, 2825--2830.
\bibitem{nasa2024} NASA Exoplanet Archive (2024). \textit{Kepler Cumulative Table}. \url{https://exoplanetarchive.ipac.caltech.edu/}
\end{thebibliography}

\end{document}